\subsection{Isotropic subspaces.}

DEFINITION 1. --- Given a module E over the ring A, an element x of E is said to be isotropic if \(\Phi(x, x) = 0\) . A submodule F of E is said to be

\begin{enumerate}
\def\labelenumi{\arabic{enumi})}
\item
  isotropic if there exists an element \(x \neq 0\) of F orthogonal to F;
\item
  totally isotropic if the restriction of \(\Phi\) to F is zero.
\end{enumerate}

When the module E is equipped with a quadratic form Q, an element of E will be said to be isotropic (resp. a submodule of E will be said to be isotropic, or totally isotropic) if this element is isotropic (resp. if this submodule is isotropic, or totally isotropic) with respect to the bilinear form associated with Q.

An isotropic vector is none other than a vector orthogonal to itself. To say that a submodule F is isotropic means that \(F \cap F^{0} \neq \{0\}\) , or equivalently that the restriction of \(\Phi\) to F is degenerate; a non-isotropic submodule G of E is therefore a submodule such that the restriction of \(\Phi\) to G is nondegenerate. For a submodule F of E to be totally isotropic, it is necessary and sufficient that one has \(F \subset F^{0}\) . If F is a totally isotropic submodule of E, the same is true of every submodule \(F'\) contained in F. The sum of a family of totally isotropic submodules that are pairwise orthogonal is a totally isotropic submodule. The set of totally isotropic submodules of E, ordered by inclusion, is evidently inductive; it follows that every totally isotropic submodule is contained in a maximal totally isotropic submodule.

PROPOSITION 1. --- Suppose that A is a field. Let F be a finite-dimensional non-isotropic subspace of E; then E is the direct sum of F and F⁰.

Indeed, since the restriction \(\Phi'\) of \(\Phi\) to F is nondegenerate by hypothesis, the map \(d_{\Phi'}\) from F to its dual F* associated on the right to \(\Phi'\) is injective, hence bijective since F and F* are two spaces of the same finite dimension. Consequently, for every \(y \in E\) , there exists one and only one element \(y_0\) of F such that one has \(\Phi(x, y) = \Phi(x, y_0)\) for all \(x \in F\) , that is to say \(y - y_0 \in F^0\) ; this proves that E is the direct sum of F and F \(^0\) .

COROLLARY. --- If F is a finite-dimensional subspace of E, and if Φ is nondegenerate, the following conditions are equivalent:

\begin{enumerate}
\def\labelenumi{\alph{enumi})}
\item
  F is non-isotropic.
\item
  \(\mathbf{F}^0\) is non-isotropic.
\item
  E is the direct sum of F and F \(^{0}\) .
\end{enumerate}

Proposition 1 indeed shows that \(a)\) implies \(c)\) , and \(c)\) implies \(a)\) and \(b)\) . Finally, if \(\mathbf{F}^0\) is non-isotropic, we have \(\mathbf{F} \cap \mathbf{F}^0 \subset \mathbf{F}^0 \cap \mathbf{F}^{00} = \{0\}\) ; hence \(\mathbf{F}\) is non-isotropic, which shows that \(b)\) implies \(a)\) .

DEFINITION 2. --- Let Q be a quadratic form on E. An element x of E is said to be singular (relative to Q) if Q(x) = 0. A submodule F of E is said to be:

\begin{enumerate}
\def\labelenumi{\arabic{enumi})}
\item
  singular if there exists an element \(x \neq 0\) of F which is singular and orthogonal to F;
\item
  totally singular if the restriction of Q to F is zero.
\end{enumerate}

The kernel of the quadratic module (E, Q) (§ 3, no. 4) consists of the singular elements of E°; for a submodule F to be singular, it is necessary and sufficient that its kernel be ≠ \{0\} . As Φ(x, y) = Q(x + y) − Q(x) − Q(y), every totally singular submodule ≠ \{0\} is singular. Since Φ(x, x) = 2Q(x), every singular vector is isotropic and every singular (resp. totally singular) submodule is isotropic (resp. totally isotropic); the converse holds if A is a field of characteristic ≠ 2. Every submodule contained in a totally singular submodule is itself totally singular. The sum of a family of totally singular submodules that are pairwise orthogonal is a totally singular submodule. The set of totally singular submodules of E, ordered by inclusion, is inductive; hence every totally singular submodule of E is contained in a maximal totally singular submodule.

